\section{Conclusion}
\label{sec:conclusion}

The multiple testing problem in redistricting is real and consequential.
Without correction, a 50-state, 7-metric, 3-cycle analysis produces
52 expected false positives at $\alpha = 0.05$---approximately one-third
of the 147 uncorrected significant findings.
The BH FDR procedure at $q = 0.10$ reduces this to at most 8.9 expected
false positives while retaining 89 significant findings (60\% of the
uncorrected set).

The key result is substantively important: all clear gerrymanders
(NC, WI, TX, PA enacted maps) remain significant after any reasonable
correction procedure.
These are genuine statistical outliers whose p-values ($< 0.004$) are
far below any multiple testing threshold.
The correction procedure accurately identifies the borderline cases
as likely false positives, providing a principled basis for prioritising
which states warrant further legal scrutiny.

Three implementation recommendations follow from this analysis:

\begin{enumerate}
\item The BH FDR procedure (at $q = 0.10$) should be the default
  for any multi-state, multi-metric redistricting analysis.
  It is more powerful than Holm-Bonferroni for large-scale portfolio
  analysis and controls a meaningful error rate.

\item The existing Holm-Bonferroni implementation in \texttt{bloc\_voting.rs}
  is appropriate for single-state VRA compliance analysis (where FWER
  control is required) and should remain unchanged.
  The new BH module in \texttt{fdr\_correction.rs} is for the L-track
  portfolio analysis only.

\item Expert reports should specify the correction procedure used and
  report both the uncorrected and corrected significance levels.
  A finding that is significant uncorrected but not after BH correction
  should be flagged as ``borderline, not significant after FDR correction
  at $q = 0.10$.''
\end{enumerate}

The Track S program's overall conclusion: redistricting inference can
be made litigation-grade by applying three existing statistical tools---
ESS correction (Component F1), calibrated p-values (S.1), and BH FDR
correction (S.4)---to the standard ensemble analysis.
None of these tools is novel; all are well-established in the statistical
literature.
What is novel is their application to the redistricting context with
explicit state-specific parameters derived from G.4's diagnostic framework.
